\documentclass[10pt,a4paper]{amsart}
\usepackage[margin=19mm]{geometry}
\usepackage{amsmath,amssymb,mathtools}
\usepackage[scr=rsfs,cal=euler]{mathalfa}
\usepackage{xfrac}
\usepackage[unicode,hidelinks,bookmarks=false]{hyperref}
\newcommand{\ie}{i.e.\ }
\newcommand{\cf}{cf.\ }
\newcommand{\resp}{respectively\ }
\newcommand{\Subsection}{\subsection}
\providecommand{\nobreakdash}{\nobreak\hbox{-}}
\setlength{\emergencystretch}{2em}
\multlinegap0pt
\usepackage[shortlabels]{enumitem}
\usepackage{amssymb}
\usepackage{tikz}
\usepackage{cancel}
\newtheorem{lemma}{Lemma}
\newcommand{\ad}{\mathrm{ad}}
\newcommand{\real}{\mathbb{R}}
\newcommand{\so}{\mathfrak{so}}
\title{Obstructions to Spin(7) Nahm transforms on tori}
\author{Spencer Whitehead}
\date{Source: arXiv:2607.28303v1 (CC BY 4.0)}
\begin{document}
\maketitle

\noindent\emph{Source and licence.} Obstructions to Spin(7) Nahm transforms on tori. Version arXiv:2607.28303v1 (CC BY 4.0), distributed by arXiv under the Creative Commons Attribution 4.0 International licence, http://creativecommons.org/licenses/by/4.0/, as recorded in the arXiv licence metadata for this exact version and shown on the abstract page https://arxiv.org/abs/2607.28303v1. Full licence notice: \texttt{licenses/arxiv-2607.28303-CC-BY-4.0.txt}.\par\smallskip

\noindent\textit{Editorial scope.} Counted unit: the article's lemma eq:pfaff (source0.tex lines 259--270). The article's complete original proof of it is delivered with it (source0.tex lines 271--325, one \texttt{proof} environment of 55 lines). No mathematics is written, repaired or re-derived: the statement and the proof are the article's own lines, in the article's own notation, and no retained line is substituted or re-flowed.  No further source-local context is needed: every cross-reference of every retained line resolves inside the two delivered spans.  Nothing was omitted from any retained span, so the package carries no omission record.  No retained line contains a citation, so the wrapper carries no bibliography and no bibliography entry.  No retained line contains a figure environment, an image-inclusion command or drawing code, so no image is delivered and the package declares no asset dependency.  Editorial formatting only, in addition: an AMS \texttt{amsart} preamble that restates the article's own declarations and macros used by the retained lines, namely the environments \texttt{lemma} on separate counters, together with the standard packages those lines need.  The retained environments print in this document as Lemma 1 in order of appearance, whereas the article prints the same items with the numbering of its own full document.  The complete native source of the article as distributed in the arXiv e-print for this exact version is bundled unchanged as \texttt{sources/206405/source0.tex}.
\begin{lemma}
    For a matrix $M$,
    \[
        \tau(c(M)) = -\frac{1}{2} \mathrm{Tr}\,M - \frac{i}{2} \langle \omega, M \rangle.
    \]
    If $|I| = 2q, q > 1$, then
    \begin{equation}
        \label{eq:pfaff}
        \tau(c(e_I)) = \sum_\sigma \mathrm{sgn}(\sigma) \prod_{j=1}^q \tau(c(e_{\sigma(2j-1)}) c(e_{\sigma(2j)}));
    \end{equation}
    the sum runs over all permutations $\sigma$ with the property that $\sigma(2j-1) < \sigma(2j)$ for each $j = 1, \ldots, q$, and $\sigma(1) < \sigma(3) < \cdots < \sigma(2q-1)$.
\end{lemma}

\begin{proof}
    Let $v$ be a unit $-4i$ eigenvector for $c(\omega)$.
    Under the adjoint map $\ad(\omega) = [\omega, \cdot]$, there is a Lie algebra decomposition
    \[
        \so(8) = \ker \ad(\omega) \oplus \mathrm{im}\, \ad(\omega).
    \]
    Let $M$ be any skew-symmetric matrix, and compute
    \begin{align*}
        \tau(c([M, \omega])) = \langle c([M, \omega])v, v \rangle = -\frac{1}{2} \langle [c(M), c(\omega)]v, v \rangle = -\frac{1}{2} (-4i \tau(c(M)) + 4i \tau(c(M))) = 0. 
    \end{align*}

    Note that $\ker \ad(\omega) = \mathfrak{u}(4) = \langle \omega \rangle \oplus \mathfrak{su}(4)$.
    For $M \in \ker \ad(\omega)$, $c(M)v = -\frac{1}{4i} c(\omega) c(M) v$; that is, $c(M) v$ is again in the $-4i$ eigenspace of $c(\omega)$.
    It follows that for $M_1, M_2 \in \ker \ad(\omega)$, $[c(M_1), c(M_2)]v = 0$.
    Thus $\tau(c(\mathfrak{su}(4))) = 0$.
    Finally, $\tau(c(\omega)) = -4i$ whence
    \[
    \tau(c(M)) = -\frac{1}{2} \mathrm{Tr}\, M -\frac{i}{2} \langle \omega, M \rangle.
    \]

    The Wick formula is proved by induction.
    Let $I'_r$ denote the sub-index of $(i_1, \ldots, i_{2q})$ in which $i_r, i_{2q}$ are omitted.
    We show that
    \[
        \tau(c(e_I)) = \sum_{r=1}^{2q-1} (-1)^{r-1} \tau(c(e_{i_r}) c(e_{i_{2q}})) \tau(c(e_{I'_r})),
    \]
    from which the claimed Equation~\eqref{eq:pfaff} follows by a standard Pfaffian counting argument.

      For a vector $v \in \real^8$, define
    \[
        c^\pm(v) := \tfrac{1}{2}\big(c(v) \pm i\, c(\omega v)\big), 
    \]
    so $c(v) = c^-(v) + c^+(v)$.
    Note that $[c(\omega), c^\pm(v)] = \pm 2i c^\pm(v)$, and therefore
    \[
        c(\omega) c^\pm(v) = c^\pm(v) ( c(\omega) \pm 2i );
    \]
    that is, $c^\pm(v)$ maps the $i\mu$-eigenspace of $c(\omega)$ into the $i(\mu \pm 2)$-eigenspace.
    In particular, $c^-(v) P_{-4i} = 0$ because $-6i$ is not an eigenvalue of $c(\omega)$.

    We have
    \[
        \{ c(e_a), c^+(e_b) \} = -\delta_{ab} - i \omega_{ab} = \tau(c(e_a) c(e_b)),
    \]
    the last equality following from the above formula for $\tau(c(E_{ab}))$ where $E_{ab}$ is the skew-symmetric matrix with a $1$ in the $(a,b)$ position, a $-1$ in the $(b,a)$ position, and zeroes elsewhere.

    Compute
    \begin{align*}
        P_{-4i} c(e_I) P_{-4i} &= P_{-4i} c(e_{i_1}) \cdots c(e_{i_{2q-1}}) c^+(e_{i_{2q}}) P_{-4i} \\
                               &= P_{-4i} \left(\sum_{r=1}^{2q-1} (-1)^{2q-1-r} \{c(e_{i_r}), c^+(e_{i_{2q}}) \} c(e_{I'_r})\right) P_{-4i} \\
                               &= \sum_{r=1}^{2q-1} (-1)^{r-1} \tau(c(e_{i_r})c(e_{i_{2q}}))P_{-4i} c(e_{I'_r}) P_{-4i} \\
                               &= \sum_{r=1}^{2q-1} (-1)^{r-1} \tau(c(e_{i_r})c(e_{i_{2q}}))\tau(c(e_{I'_r})) P_{-4i}
    \end{align*}
    as desired.
\end{proof}

\end{document}
